%Chargement des paquets
\usepackage{amsmath}
\usepackage{amsfonts}
\usepackage{amssymb}
\usepackage{enumerate}
\usepackage{mathtools}
\usepackage{bbm}
\usepackage{xparse, etoolbox}
\usepackage{enumerate}
\usepackage{enumitem}
\usepackage{mathabx}
\usepackage{babel}[french]

%environment
\usepackage{amsthm}
\usepackage{keytheorems}
\usepackage{hyperref} 

%Interface théorème
\newkeytheorem{lem}[
	name = Lemme
]

\newkeytheorem{prop}[
	name = Proposition
]

\newkeytheorem{thm}[
	name = Théorème
]

\newkeytheorem{coro}[
	name = Corollaire
]

\renewcommand*{\proofname}{Démonstration}

\theoremstyle{definition}
\newtheorem{defn}{Définition}

\theoremstyle{definition}
\newtheorem*{nota}{Notation}

\theoremstyle{definition}
\newtheorem*{limi}{Liminaire}

\theoremstyle{remark}
\newtheorem{rmq}{Remarque}

\theoremstyle{plain}
\newtheorem*{fait}{Fait}


%Ensembles de nombres
\newcommand{\N} {\mathbb{N}}
\newcommand{\Ne}{\N^\ast}
\newcommand{\Z} {\mathbb{Z}}
\newcommand{\Q} {\mathbb{Q}}
\newcommand{\R} {\mathbb{R}}
\newcommand{\Rb}{\overline{\mathbb{R}}}
\newcommand{\Rp}{\R_+}
\newcommand{\Rm}{\R_-}
\newcommand{\K} {\mathbb{K}}
\newcommand{\Ket} {\mathbb{K^{\ast}}}
\newcommand{\Cx}{\mathbb{C}}

%Ensembles, tribus
\newcommand{\A}{\mathbb{A}}
\renewcommand{\P}{\mathcal{P}}
\newcommand{\T}{\mathcal{T}}
\newcommand{\F}{\mathcal{F}}
\newcommand{\C} {\mathcal{C}}
\newcommand{\ouv}{\mathcal{O}}
\newcommand{\B}{\mathcal{B}}
\newcommand{\M}{\mathcal{M}}
\newcommand{\etag}{\mathcal{E}}
\newcommand{\etagOp}{\etag^{+}(\Omega)}
\newcommand{\etagO}{\etag(\Omega)}
%%Lp avant quotientage
\newcommand{\Lic}{\mathcal{L}^1}
\newcommand{\Lpc}{\mathcal{L}^p}
\newcommand{\Linfc}{\mathcal{L}^{\infty}}
\newcommand{\Lifull}{\Lic(\Omega, \B, m)}
\newcommand{\Lpfull}{\Lpc(\Omega, \B, m)}
\newcommand{\Linffull}{\Linfc(\Omega, \B, m)}
%%Lp après quotientage
\newcommand{\Li}{L^1}
\newcommand{\Ld}{L^2}
\newcommand{\Lp}{L^p}
\newcommand{\Linf}{L^{\infty}}
\newcommand{\Listd}{\Li(\Omega, m)} 
\newcommand{\Ldstd}{\Ld(\Omega, m)} 
\newcommand{\Lpstd}{\Lp(\Omega, m)} 

%Quelques opérateurs
\DeclareMathOperator{\codim}{codim}
\DeclareMathOperator{\rg}{rg}
\DeclareMathOperator{\tr}{tr}
\DeclareMathOperator{\vect}{Vect}
\DeclareMathOperator{\ima}{Im}
\DeclareMathOperator{\id}{Id}
\DeclareMathOperator{\sign}{sign}
\DeclareMathOperator{\ext}{ext}
\newcommand{\ind}{\mathbbm{1}}
\newcommand*{\spr}[2]{\left(\,#1\,\middle|\,#2\,\right)}
\newcommand{\interior}[1]{%
  {\kern0pt#1}^{\mathrm{o}}%
}
\DeclareMathOperator{\card}{card}
\DeclareMathOperator{\ppcm}{ppcm}

%Commandes de pure flemme
\newcommand{\veps}{\varepsilon}
\newcommand{\intab}{\int_{a}^{b}}
\newcommand{\intR}{\int_{-\infty}^{\infty}}
\newcommand{\intinf}{\int_{- \infty}^{+ \infty}}
\newcommand{\equi}{\Leftrightarrow}
\newcommand{\Kev}{$\K$-espace vectoriel }
\newcommand{\Kevs}{$\K$-espaces vectoriels }
\newcommand{\Rev}{$\R$-espace vectoriel }
\newcommand{\Revs}{$\R$-espaces vectoriels }
\newcommand{\Cev}{$\Cx$-espace vectoriel }
\newcommand{\Cevs}{$\Cx$-espaces vectoriels }
\newcommand{\evn}{(E, \norm{.})}
\newcommand{\interf}[2]{[#1,#2]}
\newcommand{\limb}{\lim\limits_}
\newcommand{\lebn}{\lambda_n}
\newcommand{\pp}{\text{~presque partout}}
\newcommand{\derivp}[2]{\frac{\partial #1}{\partial #2}}
\newcommand{\famiu}{(u_i)_{i \in I}}
\newcommand{\sev}{\text{~s.e.v.}}
\newcommand{\Mnp}{M_{n,p}(\K)}
\newcommand{\Mn}{M_{n}(\K)}
\newcommand{\tq}{\text{~t.q.~}}
\newcommand{\mq}{~montrer~que~}
\newcommand{\et}{\quad \text{et} \quad}
\newcommand{\ou}{\text{~ou~}}
\newcommand{\Kx}{\K[X]}


%Commandes de gars chelou
\renewcommand{\subset}{\subseteq}

%Commande restriction, ty duckduckgo
\def\restr#1#2{\mathchoice
              {\setbox1\hbox{${\displaystyle #1}_{\scriptstyle #2}$}
              \restraux{#1}{#2}}
              {\setbox1\hbox{${\textstyle #1}_{\scriptstyle #2}$}
              \restraux{#1}{#2}}
              {\setbox1\hbox{${\scriptstyle #1}_{\scriptscriptstyle #2}$}
              \restraux{#1}{#2}}
              {\setbox1\hbox{${\scriptscriptstyle #1}_{\scriptscriptstyle #2}$}
              \restraux{#1}{#2}}}
\def\restraux#1#2{{#1\,\smash{\vrule height .8\ht1 depth .85\dp1}}_{\,#2}} 

%Quelques normes
\DeclarePairedDelimiter\abs{\lvert}{\rvert}
\DeclarePairedDelimiter\labs{\bigl\lvert}{\bigr\rvert}
\DeclarePairedDelimiter\norm{\lVert}{\rVert}
\DeclarePairedDelimiterXPP\normi[1]{}\lVert\rVert{_1}{\ifblank{#1}{\:\cdot\:}{#1}}
\DeclarePairedDelimiterXPP\normd[1]{}\lVert\rVert{_2}{\ifblank{#1}{\:\cdot\:}{#1}}
\DeclarePairedDelimiterXPP\normp[1]{}\lVert\rVert{_p}{\ifblank{#1}{\:\cdot\:}{#1}}
\DeclarePairedDelimiterXPP\normq[1]{}\lVert\rVert{_q}{\ifblank{#1}{\:\cdot\:}{#1}}
\DeclarePairedDelimiterXPP\norminf[1]{}\lVert\rVert{_\infty}{\ifblank{#1}{\:\cdot\:}{#1}}

%Bracket en angle
\DeclarePairedDelimiter\angl{\langle}{\rangle}

%Set, parenthèses
\newcommand{\set}[1]{\{ #1 \}}
\newcommand{\bigset}[1]{\bigl\{ #1 \bigr\}}
\newcommand{\Bigset}[1]{\Bigl\{ #1 \Bigr\}}
\newcommand{\bigpar}[1]{\bigl( #1 \bigr)}
\newcommand{\Bigpar}[1]{\Bigl( #1 \Bigr)}

%Opitmisation des opérateurs de comparaison
\DeclarePairedDelimiter\tio{o (}{)}
\DeclarePairedDelimiter\gro{O (}{)}


\pagestyle{fancy}

\fancyhead[C]{Décomposition de Dunford et exponentiation}
\fancyhead[R]{}

\begin{document}

\underline{Sources :} 
\begin{enumerate}[label=(\roman*)]
\item Algèbre linéaire, Grifone - page 188. 
\item Algèbre, Rombaldi - page 633.
\end{enumerate}

\underline{Leçons :}
\begin{itemize}
\item 150 : Polynômes d’endomorphisme en dimension finie. Réduction d’un endomorphisme en dimension finie. Applications.
\item 151 : Sous-espaces stables par un endomorphisme ou une famille d’endomorphismes d’un espace vectoriel de dimension finie.
Applications.
\item 152 : Endomorphismes diagonalisables en dimension finie.
\item 155 : Exponentielle de matrices. Applications.
\item 156 : Endomorphismes trigonalisables. Endomorphismes nilpotents.

\end{itemize}

\begin{nota}
$E$ est un \Kev de dimension $n$ finie.
\end{nota}

\section{Décomposition de Dunford}

\begin{lem}
\label{lem:diagsimu}
La somme de deux endomorphismes diagonalisables qui commutent est un endomorphisme diagonalisable.
\end{lem}

\begin{proof}
Soient $f$ et $g$ deux tels endomorphismes, on note $(F_i)_{i=1}^p$ les espaces propres de $f$ et $\lambda_i$ les valeurs propres
associées. On note aussi $(G_j)_{j=1}^q$ les espaces propres de $g$ et $\mu_j$ les valeurs propres associées. 
On montre d'abord que les $G_j$ sont stables par $f$ (resp. que les $F_i$ sont stables par $g$). \\
Soient $j \in \set{1, \dots, q}$ et $x \in G_j$, on a $g \circ f(x) = f \circ g (x) = f(\mu_j x) = \mu_j f(x)$ donc $f(x) \in G_j$. \\

On pose :
\[ \forall i \in \set{1, \dots, p}, \forall j \in \set {1, \dots, q} H_{i,j} = F_i \cap G_j . \]
On remarquera d'abord que, pour $i$ fixé, il est impossible que les $H_{i,j}$ soient tous réduits à $\set{0}$,
en effet, on a $E = \bigoplus_{j=1}^q G_j$ et $F_i \subset E$. 
On va montrer que :
\[ \forall i \in \set{1, \dots, p }, F_i = \bigoplus_{j=1}^q H_{i,j} . \]
Soit $i$ fixé et $x \in F_i$, on écrit $x = g_1 + \dots + g_q$ avec $g_j \in G_j, j \in \set{1, \dots, q}$. Cette décomposition est par ailleurs
unique. Or $f(x_i) = \lambda_i x_i = f(g_1) + \dots + f(g_q)$ avec $f(g_j) \in G_j, j \in \set{1, \dots, q}$ selon le point démontré plus haut.
Par unicité de la décomposition, on conclut que :
\[ \forall j \in \set{1, \dots, q}, f(g_j) = \lambda_i g_j \text{ soit } g_j \in F_i \cap G_j = H_{i,j} , \]
et cette décomposition étant unique, on a le résultat cherché. \newline

Ce résultat étant par ailleurs vrai pour tout $i$, on peut écrire :
\[ E =  \bigoplus_{i, j} H_{i,j} \]
et les matrices des endomorphismes $f$ et $g$ sont clairement diagonales sur une base composée des bases des $H_{i,j}$.
\end{proof}

\begin{lem}
\label{lem:sommenil}
La somme de deux endormophismes nilpotents qui commutent est un endomorphisme nilpotent.
\end{lem}

\begin{proof}
Soient $n$ et $n'$ deux tels endomorphismes d'indices de nilpotences respectifs $r$ et $s$. Alors :
\begin{align*}
(n+n')^{r+s} & = \sum_{k=0}^{r+s} \binom{k}{r+s} n^k n'^{r+s-k} \\
& = \sum_{k=0}^{r} \binom{k}{r+s} n^k n'^{r+s-k} + \sum_{k=r+1}^{r+s} \binom{k}{r+s} n^k n'^{r+s-k}
& = 0 + 0 
\end{align*}
\end{proof}

\begin{thm}[Décomposition de Dunford]
Soit $f$ un endomorphisme de $E$ dont le polynôme caractéristique est scindé. Alors, $f$ se décompose de manière unique sous la forme :
\[ f = d + n , \]
où $d$ est un endomorphisme diagonalisable et $n$ un endomorphisme nilpotent tels que $d$ et $n$ commutent.
\end{thm}

\begin{proof}
\begin{enumerate}
\item Existence \\

Comme le polynôme caractéristique de $f$ est scindé, selon le le lemme des noyaux, on a :
\[ E = N_{\lambda_1} \oplus \dots \oplus N_{\lambda_p} , p \in \N \]
avec les $N_{\lambda_i}$ les sous espaces caractéristiques de $f$. 
On peut donc trouver une base $B$ de $f$ dans laquelle la matrice de $f$ est :
\begin{gather*} A = M_B(f) = \begin{pmatrix} 
\begin{pmatrix} 
\lambda_1 & & * \\
 & \ddots & \\
0 & & \lambda_1 \\
\end{pmatrix} & & & 0 \\
& \begin{pmatrix} 
\lambda_2 & & * \\
 & \ddots & \\
0 & & \lambda_2 \\
\end{pmatrix} & & \\
 & & \ddots & \\
0 & & & \begin{pmatrix} 
\lambda_p & & * \\
 & \ddots & \\
0 & & \lambda_p \\
\end{pmatrix} 
\end{pmatrix} \\
= \begin{pmatrix} 
\begin{pmatrix} 
\lambda_1 & & 0 \\
 & \ddots & \\
0 & & \lambda_1 \\
\end{pmatrix} & 0 \\
 & \ddots & \\
0 & & \begin{pmatrix} 
\lambda_p & & 0 \\
 & \ddots & \\
0 & & \lambda_p \\
\end{pmatrix} 
\end{pmatrix} 
+ \begin{pmatrix} 
\begin{pmatrix} 
0 & & * \\
 & \ddots & \\
0 & & 0 \\
\end{pmatrix} & 0 \\
 & \ddots & \\
0 & & \begin{pmatrix} 
0 & & * \\
 & \ddots & \\
0 & & 0 \\
\end{pmatrix} 
\end{pmatrix} 
\end{gather*}

En notant $D$ la première matrice et $N$ la seconde, on a $A = D + N$ et soient $d$ et $n$ les endomorphismes de $E$ qui, 
dans la base $B$ sont représentés par $D$ et $N$ respectivement. 
ses endomorphises sont respectivement diagonalisable et nilpotent et ils commutent 
(les matrices $D$ et $N$ commutent, on le démontre en effectuant un produit par bloc).
On a donc démontré l'existence.

\item Unicité : on considère une seconde telle décomposition de $f$ soit $f = d + n = d' + n'$, on a alors $h = d - d' = n' - n$. 
On va montrer que $d$ et $d'$ commutent et que $n$ et $n'$ commutent. \\

On remarque d'abord que $d'$ commute avec $f$ (car il commute avec $d'+n'$). 
On va montre que les sous-espaces caractéristiques de $f$ sont stables par $d'$. \\
Prenons $x \in N_{\lambda_i}$, puisque $d'$ et $f$ commutent on a :
\[ (f-\lambda_i)^{\alpha_i} \circ d'(x) = d' \circ (f-\lambda_i)^{\alpha_i} (x) = 0 , \]
donc $d'(x) \in N_{\lambda_i}$. \\

Soit $x \in E, x= x_1 + \dots + x_p$, avec $x_i \in N_{\lambda_i}$.
Ainsi, en notant $q_i$ le projecteur sur $N_{\lambda_i}$ on a : 
\begin{gather*}
d' \circ q_i(x) = d'(x_i) \qquad \text{ et} \\
q_i \circ d'(x) = q_i \circ d'(x_1) + \dots + q_i \circ d'(x_p) = q_i \circ d'(x_i) = d'(x_i)
\end{gather*}
Donc $d'$ commute avec chaque projecteur $q_i$. \\

De plus, on a vu que la matrice de $d$ était diagonale dans une base $B$ constituée des bases des sous-espaces caractéristiques.
On peut donc écrire $d = \lambda_1 q_1 + \dots \lambda_p q_p$ et conclure que $d$ commute avec $d'$. \\

On démontre de la même manière que $n'$ commute avec $f$ et $d$. Enfin :
\[ n' \circ n = n' \circ (f-d) = (f-d) \circ n' = n \circ n' . \]

Donc, selon les lemmes $\ref{lem:diagsimu}$ et $\ref{lem:sommenil}$ l'endormophisme $h$ est un endomorphisme diagonalisable ET
nilpotent, c'est donc l'endomorphisme nul, ce qui conclut la preuve d'unicité.
\end{enumerate}

\end{proof}

\section{Exponentielle d'un endomorphisme}

On suppose que $\K = \C$, et que $\norm{.}$ est une norme sur$\L(E)$. On considère $u \in \L(E)$.

\begin{lem}
En notant $u = d+v$ le décomposition de Dunford de $u$, on a :
\[ e^u = e^d e^v =e^d \sum_{k=0}^{q-1} \dfrac{v^k}{k!} , \]
où $q \geq 1$ est l'indice de nilpotence de $v$ et $e^u$ est un polynôme en $u$.
\end{lem}

\begin{proof}
Comme $d$ et $v$ commutent, on a :
\begin{align*}
e^{u+v} & = \sum_{n=0}^{\infty} \dfrac1{n!} (u+v)^n \\
& = \sum_{n=0}^{\infty} \dfrac1{n!} \sum_{k=0}^n \binom{n}{k} u^k v^{n-k} \\
& = \sum_{n=0}^{\infty} \sum_{k=0}^n \dfrac{u^k}{k!} \dfrac{v^{n-k}}{(n-k)!} \\
& = e^u e^v .
\end{align*}
La simplifciation $e^v =  \sum_{k=0}^{q-1} \dfrac{v^k}{k!}$ s'obtient en remarquant que $v^k = 0$ pour $k \geq q$.
\end{proof} 

\begin{thm}
Si $u=d+v$ est la décomposition de Dunford de $u$, alors celle de $e^u$ est donnée par 
\[ e^u = e^d +e^d(e^v-\id), \]
avec $e^d$ diagonalisable et $e^d(e^v-\id)$ nilpotente.
\end{thm}

\begin{proof}
On a :
\begin{align*}
e^u & = e^d e^v = e^d \sum_{k=0}^{q-1} \dfrac{v^k}{k!} = e^d \left ( \id + v \sum_{k=1}^{q-1} \dfrac{v^{k-1}}{k!} \right ) \\
& = e^d + v e^d \sum_{k=1}^{q-1} \dfrac{v^{k-1}}{k!} = e^d + v \cdot w
\end{align*}
en remarquant que $v$ et $e^d$ commutent (car $v$ et $d$ commutent). De même, $v$ et $w$ commutent, on a donc 
$(v \cdot w)^q = v^q w^q = 0$ en notant $q$ l'indice de nilpotence de $v$. De plus, $e^d$ étant diagonale, on a obtenu, par 
unicité de ce suivant, la décomposition de Dunford de $e^u$. La partie nilpotente de cette décomposition s'écrit aussi :
\[ v \cdot w =  e^d \sum_{k=1}^{q-1} \dfrac{v^k}{k!} = e^d(e^v-\id) . \]
\end{proof}

\begin{coro}
Un endomorphisme $u$ est diagonalisable si, et seulement si, $e^u$ est diagonalisable.
\end{coro}

\begin{proof}
On écrit $u = d+v$ la décomposition de Dunford de $u$. Si $u$ est diagonalisable, alors $v=0$ et, par suite, $e^d(e^v-\id) = 0$ donc
$e^u$ est diagonalisable. \\
Réciproquement, si $e^d(e^v-\id) = 0$, alors $e^v = \id$ puisque $e^d$ est inversible. On a donc :
\[ \sum_{k=0}^q \dfrac{v^k}{k!} = \id \iff \sum_{k=1}^q \dfrac{v^k}{k!} = 0 , \]
avec $q$ l'indice de nilpotence de $v$. Autrement dit, $P(X) = \sum_{k=1}^q \dfrac{X^k}{k!}$ est un polynôme annulateur de $v$.
Or $X^q$ est le polynôme minimal de $v$, ce qui impose $q=1$ (par identification du monôme de degré $q$), soit $v=0$ et donc
$u$ diagonalisable. 
\end{proof}




\end{document}